\begin{definition}[Furstenberg index]\label{deffurindex}
    Suppose $(s,t;n,k)$ is admissible, we define $\bF(s,t;n,k)$ as follows.
    
    (a) If $s=0$, define
    \begin{equation}
        \bF(0,t;n,k)= \max\{0,t-k(n-k)\}.
    \end{equation}

    (b) If $s>0$, write $s=d+\si$ where $d\in\N, \si\in(0,1]$. If also $t\in [0,(k-d-1)(n-k)]$, define
    \begin{equation}
        \bF(s,t;n,k)= s.
    \end{equation}


    If we are not in Case (a) and (b), then $s\in (0,k]$ and $t\in ((k-d-1)(n-k),(k+1)(n-k)]$,
    we write 
\begin{equation}\label{expression}
    \begin{cases}
        s=d+\si\\
        t=(k-d-1)(n-k)+(d+2)m+\tau,
    \end{cases}
\end{equation} 
    where $d\in\{0,\dots,k-1\}$, $\si\in (0,1]$,  $m\in\{0,\dots,n-k-1\}$, $\tau\in (0,d+2]$. 

(c) If $\tau\in(0,2]$, define
\begin{equation}
\begin{split}
    \bF(s,t;n,k)&=s+m+\min\{\tau,\frac12(\si+\tau),1\}\\
    &= d+m+\min\{\si+\tau,\frac32\si+\frac12\tau,\si+1\}.
\end{split}
\end{equation}

(d) If $\tau\in (2,d+2]$, define
\begin{equation}
\begin{split}
    \bF(s,t;n,k)&=s+m+\min\{\tau,\frac12(\si+\tau),1\}\\
    &= s+m+1.
\end{split}
\end{equation}
    
    
\end{definition}

\begin{remark}
    {\rm In (c) and (d), $\bF$ has the same expression, but we separate them for the convenience of later discussions.
    
    Through this definition, we can calculate $\bF(s,t;2,1)=\min\{s+t,\frac32s+\frac12t,s+1\}$ for $s\in (0,1], t\in [0,2]$. We see that Conjecture \ref{conj2d} is equivalent to saying that $\inf_{E:\ (s,t;2,1)\textup{-set}} \#E\gtrsim_\e p^{\bF(s,t;2,1)-\e}$. }
\end{remark}

The following is a construction of a $(s,t;2,1)$-set, which provides evidence on why Conjecture \ref{conj2d} should be reasonable.

\begin{proposition}\label{prop2d}
For $0\le t\le 2$, there exists a $(0,t;2,1)$-set $E$ such that
    \begin{equation}
        \#E\lesssim p^{\max\{0,t-1\}}.
    \end{equation}
For $0<s\le 1, 0\le t\le 2$.\, there exists a $(s,t;2,1)$-set $E$ such that    
    \begin{equation}\label{f21}
        \#E\lesssim p^{\min\{s+t,\frac32s+\frac12t,s+1\}}.
    \end{equation}
\end{proposition}

\begin{proof}
When $s=0$, we construct $E=\bigcup_{L\in \cL} Y(L)$ as follows. If $t\le 1$, choose $\cL$ to be a set of $p^t$ lines passing through $(0,0)$ and $Y(L)=(0,0)$. Then $\#E=1=p^0$. If $t>1$, we choose $p^{t-1}$ points $E$ on $\Fp\times \{0\}$, and for each $x\in E$, we choose $p-1$ lines through $x$ that are not parallel to $\Fp\times \{0\}$. If the line $L$ is through $x\in E$, then we let $Y(L)=x$. We see that $E$ is a $(0,t;2,1)$-set with cardinality $p^{t-1}$. (Here, $p^{t-1}$ should actually be the integer $\lceil p^{t-1}\rceil$, but we still write it as $p^{t-1}$ for simplicity, thinking of it as an integer. Same convention applies in the whole paper.) 

Consider the case when $s>0$. There are three terms in the ``min" on the RHS of \eqref{f21}. We note that: when $t\le s$, minimum achieves at $s+t$; when $s\le t\le 2-s$, minimum achieves at $\frac32s+\frac12t$; when $2-s\le t$, minimum achieves at $s+1$.
\begin{enumerate}
    \item If $E=\bigcup_{L\in\cL} Y(L)$ is a $(s,t;2,1)$-set, then $\#E\le \sum_{L\in\cL} \#Y(L)\le p^{s+t}$.

    \item Choose $\cL$ to be a set of $\frac12 p^t$ lines that are not parallel to the line $\Fp\times\{0\}$. Let $Y(L)=L\cap (\Fp\times \{1,2,\dots,p^s\})$. Then $E$ is a $(s,t;2,1)$-set with $E=\bigcup_{L\in\cL} Y(L)\subset \Fp\times \{1,\dots,p^s\}$. Therefore, $\#E\le p^{s+1}$.

    \item It remains to consider the case when $s+t\le 2, s\le t$. This example is of Szemer\'edi-Trotter type. The heuristic is shown in Figure \ref{F0}, where we have $\sim p^{\frac{t-s}{2}}$ directions and for each direction, there are $\sim p^{\frac{t+s}{2}}$ lines. We also have a rectangle of size $\sim p^{s}\times p^{\frac{s+t}{2}}$. The rectangle intersects each line $L$ at $\sim p^s$ points, which is defined to be $Y(L)$. We also see $E=\bigcup_L Y(L)$ is contained in the rectangle. Therefore, $\#E\lesssim p^{\frac32s+\frac12t}$.

    
    Let $\cL=\{ y=ax+b: a=1, \dots,  p^{\frac{t-s}{2}},b=1, \dots,  p^{\frac{s+t}{2}} \}$. Here, we view $a,b$ as elements in $\Fp$, and hence $y=ax+b$ as a line in $\Fp^2$. For each $L\in \cL$, let $Y(L)=\{(x,y)\in L:x=1,\dots, \frac12p^s\}$. We see that $E=\bigcup_{L\in \cL}Y(L)$ is contained in $\{(x,y): x=1,\dots,  \frac12p^s, y=1,\dots, 2 p^{\frac{s+t}{2}} \}$. Therefore, $\#E\lesssim p^{\frac32s+\frac12t}$. 
\end{enumerate}
\end{proof}

\begin{figure}[ht]
\centering
\begin{minipage}[b]{0.45\linewidth}
\includegraphics[width=5cm]{F0.png}
\caption{}
\label{F0}
\end{minipage}
\end{figure}

\bigskip
